\section*{Do wrong Sudoku branches contain different useful digits?}
For a communication mechanism to help, the sending branch must contribute something the receiving branch lacks. Sudoku offers a simple first test because the same puzzle has one target grid. At a fixed blank the archived model proposed two digits. Each proposal edited the working state and ran the native model twice. Even when both predicted grids were wrong, one might supply correct digits where the other was wrong.
\par
\textbf{Which boards qualify?} In the historical 256-board development cohort, 94 boards had two wrong branch endpoints. On 77 the endpoints differed. On 76 of those, none of the candidates from the matched serial continuation was exact. These 76 boards were chosen from archived GH200 correctness flags before inspecting per-cell results. The archive did not store the full digit grids. We recovered the native checkpoint and the learned edit, checked their SHA-256 hashes and replayed the 76 boards on CPU from the saved depth-two development states. Reconstruct the depth-12 state and rerun both conditional paths and four more unconditioned serial steps. CPU and GPU numerical differences can change an ambiguous fork; we report replay eligibility separately.
\par
\textbf{How to read the ceiling.} Let $a_i,b_i$ be the two endpoint digit predictions and $y_i$ the true digit at cell $i$. A label-aware chooser can assemble the exact solution from those two grids if and only if every cell satisfies $a_i=y_i$ or $b_i=y_i$. A single shared error disproves sufficiency for that board. Reciprocal contributions require at least one cell correct only in A and at least one correct only in B. These are retrospective diagnostics; labels never choose digits in a deployed model. Averaging the endpoint logits supplies a separate label-free check.
\par
\textbf{CPU replay.} Of the 76 archived cases, 50 still had two different, wrong endpoints and a serial miss on CPU. The selected blank and proposed pair of digits matched the archived values on 1/76 boards. On 14/50 replay-eligible boards, each branch was correct on at least one cell missed by the other. A chooser with access to the true digits could assemble 0/50 exact boards. The median number of cells wrong in both endpoints was 30. Averaging the endpoint logits solved 0/50 boards. These counts apply to the selected failures. They do not estimate full-cohort improvement.
Of all 76, averaging solved five, but each of those already had an exact CPU branch endpoint. Agreement on only one selected blank-and-digit pair prevents attribution to the historical GH200 trajectories. Numerical divergence is a possible explanation; the cause has not been established.
\par
\textbf{Verified on the original runtime.} Slurm job 8930277 reproduced all 76 fingerprints, including the chosen cell, proposed digits, four branch digit readouts and stepwise exact flags. Each branch supplied at least one correct digit missed by its peer on 17/76 boards. Across the 76 failures the median shared-error count was 29. No target could be assembled from the two grids, and average logits solved none. These fingerprints are stronger than matching aggregate totals, but the archive has no original complete grids for a bitwise comparison. The job completed on one GH200 in 32 seconds (0.0089 GPU-hours). The full new predictions are retained on project storage.

\par
\textbf{Registered follow-up.} On fresh original puzzles, freeze both branches and their continuation. At a fixed step deliver the peer's correction to corresponding disagreement cells. Compare with the same vectors permuted among the same cells, and with the untouched parent. Equalize observations, total block calls, and post-cast correction size. Choose sites without labels. Report paired exact gains and losses over all puzzles and inspect which later computation changes. The independent reviewer emphasized that damaging the shuffled arm is insufficient: aligned communication must improve the unmodified trajectory too. Compare a serial solver and an online whole-trajectory selector before training. The following page reports the outcome.
